\chapter{Introduction}

Automation is influencing almost every part of daily life, and transportation is one of the areas changing most rapidly. 
\acp{AV} are at the forefront of this change, revolutionizing the transportation sector. Their development has accelerated because of recent progress in \ac{AI}, 
sensing systems, computing, and communication technologies. These innovations could enhance traffic efficiency, increase travel comfort, improve road
safety, and, most importantly, reshape future mobility.

\par However, fully autonomous traffic systems do not yet exist anywhere in the world. \acp{AV} will have to operate in mixed traffic together with 
\acp{HDV}. This creates important challenges because different road users provide different types and amounts of information. 
While \acp{AV} can rely on
rich on-board sensing and may also use \ac{V2X} communication \cite{abedi2025precisehdvpositioningsafetyaware}, 
the intentions of human drivers are not directly observable and must be inferred from their behavior in traffic \cite{Schwarting2018}. 
Consequently, their intentions and decisions cannot be measured directly,  
but instead have to be estimated from observable cues such as vehicle motion, lane changes, relative positions, and interactions with nearby vehicles.

\par For this reason, human driver behavior prediction is an important problem in autonomous driving. If an \ac{AV} can estimate what nearby human drivers are likely to do next, 
it can make safer and better-informed decisions. Anticipating likely braking, lane-changing, or car-following behavior can support downstream tasks such as risk 
assessment, trajectory prediction, and motion planning \cite{Lefevre2014,Madjid_2026,Schwarting2018}.

\par Over the last two decades, researchers have developed three broad families of behavioral
models for road users: rule-based models, probabilistic models, and learning-based methods \cite{10054035,Rudenko_2020,Lefevre2014}.
In parallel, several game-theoretic decision-making paradigms for \acp{AV} have emerged \cite{automotiveinnovation}. Rule-based models are simple and easy to understand, 
but they are often too rigid to capture the variability of real human behavior. Learning-based methods can learn complex patterns from large datasets, but many of them 
are difficult to interpret and are considered as black boxes. 
Probabilistic models provide a useful middle ground because they can explicitly represent uncertainty, capture temporal dependencies, and still remain interpretable.

\par This thesis focuses on probabilistic modeling of high-level human driving behavior in highway traffic. 
Rather than building the full control or planning stack of an \ac{AV}, this work addresses an important upstream problem: predicting the likely future maneuvers of surrounding human drivers. 
The goal is to develop a model that can estimate likely future maneuvers while also representing uncertainty in a clear and interpretable way.
In this thesis, future maneuvers are represented internally as discrete action states.

\par To achieve this, the thesis uses a \ac{DBN} to model the temporal relationship between observable driving features and hidden behavioral states. 
Since driver intentions are not directly observable, the model represents them through latent variables and infers them probabilistically from trajectory-based observations \cite{murphy2002dynamic}. 
In particular, the \ac{DBN} is designed to capture latent factors such as driving context and action dynamics, and to produce probabilistic predictions of future high-level maneuvers.
The model is trained and evaluated using the highD dataset, a naturalistic highway driving dataset containing detailed vehicle trajectory recordings \cite{krajewski2018highddatasetdronedataset}. From these trajectories, 
relevant features are extracted to describe motion, surrounding traffic context, and interaction behavior. These features serve as observations for the \ac{DBN}.

\par The thesis studies whether such a probabilistic framework can provide accurate, interpretable, and computationally efficient short-horizon predictions of future human driving maneuvers. 
In particular, it examines the suitability of the dataset, the structure of the \ac{DBN}, the usefulness of the selected observation features, the quality of future maneuver prediction, 
the interpretability of the learned states, and the runtime performance of the model for online use.

